% Fichier généré par generer_figures.py : ne pas modifier à la main.
\newcommand{\sourireiv}[1]{\draw[#1] plot coordinates {(92.0,28.66) (92.5,28.39) (93.0,28.13) (93.5,27.86) (94.0,27.59) (94.5,27.33) (95.0,27.06) (95.5,26.80) (96.0,26.54) (96.5,26.28) (97.0,26.01) (97.5,25.75) (98.0,25.49) (98.5,25.23) (99.0,24.97) (99.5,24.71) (100.0,24.46) (100.5,24.20) (101.0,23.94) (101.5,23.68) (102.0,23.43) (102.5,23.17) (103.0,22.91) (103.5,22.66) (104.0,22.40) (104.5,22.15) (105.0,21.89) (105.5,21.64) (106.0,21.38) (106.5,21.13) (107.0,20.87) (107.5,20.62) (108.0,20.36)};}
\newcommand{\sourirevl}[1]{\draw[#1] plot coordinates {(92.0,34.21) (92.5,33.59) (93.0,32.97) (93.5,32.36) (94.0,31.76) (94.5,31.17) (95.0,30.58) (95.5,30.00) (96.0,29.42) (96.5,28.86) (97.0,28.29) (97.5,27.74) (98.0,27.19) (98.5,26.65) (99.0,26.11) (99.5,25.58) (100.0,25.05) (100.5,24.54) (101.0,24.02) (101.5,23.51) (102.0,23.01) (102.5,22.52) (103.0,22.03) (103.5,21.54) (104.0,21.06) (104.5,20.59) (105.0,20.12) (105.5,19.65) (106.0,19.20) (106.5,18.74) (107.0,18.29) (107.5,17.85) (108.0,17.41)};}
\newcommand{\pointsbloomberg}[1]{\foreach \p in {(90.0,29.62),(95.0,27.33),(100.0,24.52),(105.0,21.42),(110.0,19.60)} {\fill[#1] \p circle (1.6pt);}}
